\section{Ordinamento in Tempo Lineare}

\subsection{Definizione}
In questa lezione si distinguono:
\begin{itemize}
	\item algoritmi di ordinamento \textbf{basati sul confronto},
	\item algoritmi \textbf{non} basati sul confronto (CountingSort, BucketSort).
\end{itemize}

Per gli algoritmi basati sul confronto vale il limite inferiore nel caso peggiore:
\[
\Omega(n\log n).
\]

\subsection{Intuizione}
Il limite inferiore discende dagli alberi di decisione: con $n$ elementi distinti ci sono $n!$ permutazioni possibili, quindi l'altezza dell'albero deve soddisfare $2^h\ge n!$.

\subsection{Motivazione}
Per valori numerici con struttura nota (intervallo limitato, distribuzione uniforme) si possono superare i limiti degli ordinamenti per confronto nel modello generale.

\subsection{Algoritmo}
\paragraph{CountingSort (slide 12--16).}
\begin{algorithm}[H]
	\caption{\textsc{CountingSort}$(A,k)$}
	\KwIn{$A$ array di interi in $[0,k]$}
	\KwOut{$B$ array ordinato}
	$B[1:n]$ nuovo array\;
	$C[0:k]$ nuovo array inizializzato a 0\;
	\For{$j\gets1$ \KwTo $A.length$}{
		$C[A[j]]\gets C[A[j]]+1$\;
	}
	\For{$i\gets1$ \KwTo $k$}{
		$C[i]\gets C[i]+C[i-1]$\;
	}
	\For{$j\gets A.length$ \KwTo $1$}{
		$B[C[A[j]]]\gets A[j]$\;
		$C[A[j]]\gets C[A[j]]-1$\;
	}
	\Return $B$\;
\end{algorithm}

\paragraph{BucketSort (slide 17--28).}
\begin{algorithm}[H]
	\caption{\textsc{BucketSort}$(A)$}
	\KwIn{$A$ array di decimali positivi in $[0,1)$}
	\KwOut{$C$ array ordinato}
	$n\gets A.length$\;
	$B[0:n-1]$ array di liste vuote\;
	$C[1:n]$ array vuoto\;
	\For{$i\gets1$ \KwTo $n$}{
		$B[\lfloor n\cdot A[i]\rfloor].add(A[i])$\;
	}
	\For{$i\gets0$ \KwTo $n-1$}{
		\textsc{InsertionSort}$(B[i])$\;
		Aggiungi a $C$ tutti i valori ordinati di $B[i]$\;
	}
	\Return $C$\;
\end{algorithm}

\subsection{Esempio guidato}
\paragraph{CountingSort.}
Con $A=[2,5,3,0,2,3,0,3]$, $k=5$:
\begin{itemize}
	\item conteggi: $C=[2,0,2,3,0,1]$,
	\item cumulati: $C=[2,2,4,7,7,8]$,
	\item output finale: $B=[0,0,2,2,3,3,3,5]$.
\end{itemize}

\paragraph{BucketSort.}
Con $A=[0.78,0.17,0.39,0.26,0.72,0.94,0.21,0.12,0.23,0.68]$ e $n=10$:
\begin{itemize}
	\item bucket 1: $\{0.17,0.12\}$,
	\item bucket 2: $\{0.26,0.21,0.23\}$,
	\item bucket 7: $\{0.78,0.72\}$,
	\item bucket 9: $\{0.94\}$,
	\item concatenazione finale dei bucket ordinati.
\end{itemize}

\subsection{Grafico o diagramma}
\begin{center}
\begin{verbatim}
CountingSort: A=[2,5,3,0,2,3,0,3]
C (conteggi): [2,0,2,3,0,1]
C (cumulati): [2,2,4,7,7,8]
B finale:     [0,0,2,2,3,3,3,5]
\end{verbatim}
\end{center}

\begin{center}
\begin{verbatim}
BucketSort (n=10):
[0.0,0.1):  -
[0.1,0.2):  0.17,0.12
[0.2,0.3):  0.26,0.21,0.23
[0.3,0.4):  0.39
...
[0.9,1.0):  0.94
\end{verbatim}
\end{center}

\subsection{Dimostrazione}
\paragraph{Lower bound per confronto.}
Se $l$ è il numero di foglie dell'albero di decisione e $h$ l'altezza:
\[
n!\le l\le 2^h\;\Rightarrow\;h\ge\log_2(n!)=\Omega(n\log n).
\]

\paragraph{Correttezza BucketSort (slide 21).}
Per due elementi $A[i]\le A[j]$:
\[
\lfloor nA[i]\rfloor\le\lfloor nA[j]\rfloor.
\]
Quindi $A[i]$ finisce nello stesso bucket di $A[j]$ o in uno precedente.
Se nello stesso bucket, \textsc{InsertionSort} li ordina; se in bucket diversi, la concatenazione per indice crescente mantiene l'ordine globale.

\subsection{Analisi}
\paragraph{CountingSort.}
\[
T(n)=\Theta(n)+\Theta(k)+\Theta(n)=\Theta(n+k).
\]
Se $k=O(n)$, allora $T(n)=\Theta(n)$.

\paragraph{Stabilità di CountingSort.}
La stabilità deriva dalla scansione finale da destra verso sinistra.

\paragraph{BucketSort.}
Costo generale:
\[
T(n)=\Theta(n)+\sum_{i=0}^{n-1}O(n_i^2),
\]
dove $n_i$ è la dimensione del bucket $i$.

Con input uniforme in $[0,1)$:
\[
E[T(n)] = \Theta(n).
\]

Dettaglio chiave delle slide:
\[
E[n_i^2]=2-\frac{1}{n}=O(1)
\]
quindi
\[
E[T(n)]=\Theta(n)+n\,O(1)=\Theta(n).
\]

\subsection{Errori tipici d'esame}
\begin{hardnote}
Applicare CountingSort a chiavi non intere o fuori dall'intervallo $[0,k]$ senza trasformazione valida.
\end{hardnote}

\begin{hardnote}
Dimenticare l'ipotesi di distribuzione uniforme quando si dichiara $\Theta(n)$ per BucketSort.
\end{hardnote}

\subsection{Possibili domande d'esame}
\begin{itemize}
	\item Dimostrare il lower bound $\Omega(n\log n)$ tramite alberi di decisione.
	\item Spiegare perché CountingSort è stabile.
	\item Derivare $E[T(n)]$ di BucketSort nel caso uniforme.
	\item Proporre una variante di BucketSort per decimali positivi non limitati a $[0,1)$.
\end{itemize}

% ----------------------------------------------------------------
\subsection{Contenuto richiesto ma non nel materiale ufficiale}

\subsubsection*{RadixSort}
\begin{hardnote}
Contenuto non presente nel materiale ufficiale del corso.
\end{hardnote}
